\documentclass[a4paper,11pt]{article}
\usepackage[utf8]{inputenc}
\usepackage[T1]{fontenc}
\usepackage[french]{babel}
\usepackage[left=2cm,right=2cm,top=2cm,bottom=2cm]{geometry}
\usepackage{xcolor}
\usepackage{hyperref}
\usepackage{fontawesome5}
\usepackage{parskip}

% --- Couleurs ---
\definecolor{primary}{RGB}{35, 55, 123}

\begin{document}

% --- EN-TÊTE ---
\begin{center}
    {\Large \textbf{Loïc Aron MBASSI EWOLO}}\\[4pt]
    \small
    \faMapMarker\ Cergy, France \quad
    \faPhone\ (+33) 7 59 52 33 98 \quad
    \faEnvelope\ \href{mailto:loicaron.mbassiewolo@ensea.fr}{loicaron.mbassiewolo@ensea.fr}
\end{center}
\vspace{0.5cm}

% --- DESTINATAIRE ---
\hfill
\begin{minipage}[t]{0.45\textwidth}
    \textbf{Framatome}\\
    Service Recrutement\\
    Grenoble
\end{minipage}

\vspace{1cm}

\textbf{Objet :} Candidature Stage Ingénieur Structuration de Modèles et Simulation (Avril 2026)

\vspace{0.5cm}

Madame, Monsieur,

L'ingénierie des systèmes complexes et la modélisation sont des domaines qui me passionnent depuis mes premiers cours d'UML à Polytechnique Yaoundé. Votre offre de stage sur la structuration de modèles dans l'environnement ALICES, au cœur des systèmes I\&C de centrales nucléaires, représente exactement le type de défi technique que je recherche. Actuellement en double diplôme à l'\textbf{ENSEA} et à l'\textbf{École Polytechnique de Yaoundé}, je souhaite contribuer à vos travaux en MBSE.

Mon approche de la modélisation est avant tout pratique. J'ai développé un outil en Python qui parse des fichiers XML issus de Draw.io pour générer du code backend Laravel. Ce projet m'a confronté aux subtilités du \textbf{parsing de diagrammes} et à l'automatisation de la génération de code — des compétences que je pourrais mobiliser pour l'allocation de modules fonctionnels dans votre environnement de simulation.

En parallèle, j'ai acquis une expérience en programmation système avec mon \textbf{projet TAXI} en C : simulation multi-processus d'une flotte de véhicules avec mémoire partagée et communication inter-processus. Cette approche modulaire, où chaque composant fonctionne comme une brique autonome, me semble proche de la logique que vous décrivez pour vos diagrammes d'architecture.

Mon passage au \textbf{Mila} comme Assistant de Recherche m'a appris la rigueur nécessaire pour documenter et tester des systèmes complexes. C'est cette même rigueur que j'aimerais apporter à vos développements de cas de tests et à la validation de modèles comportementaux.

Je suis disponible dès \textbf{avril 2026} pour un stage de 4 à 6 mois. Je serais ravi d'échanger sur vos projets et de vous montrer mes réalisations lors d'un entretien.

Je vous prie d'agréer, Madame, Monsieur, l'expression de mes salutations distinguées.

\vspace{1cm}

\textbf{Loïc Aron MBASSI EWOLO}\\
\textit{Élève-Ingénieur ENSEA / Polytechnique Yaoundé}

\end{document}
